% =============================================================================
% 00_preambulo.tex  |  Tema e macros dos slides
% Workshop: Técnicas de Avaliação de Políticas Públicas (EACH/USP, 2026)
% Paleta inspirada no cartaz do Seminário Discente 2026 (laranja + índigo).
% =============================================================================
\usepackage[utf8]{inputenc}
\usepackage[T1]{fontenc}
\usepackage[brazil]{babel}
\usepackage{lmodern}
\usepackage{graphicx}
\usepackage{booktabs}
\usepackage{tikz}
\usetikzlibrary{arrows.meta,positioning,shapes.geometric,calc,fit,backgrounds,decorations.pathreplacing}
\usepackage{listings}
\usepackage{xcolor}
\usepackage{tabularx}
\usepackage{array}

\graphicspath{{../figuras/}}

% ---- Cores --------------------------------------------------------------------
\definecolor{laranja}{HTML}{F26B21}
\definecolor{indigo}{HTML}{4A55A2}
\definecolor{cinza}{HTML}{8E93B0}
\definecolor{verde}{HTML}{2E8B57}
\definecolor{fundo}{HTML}{F8F8F8}
\definecolor{indigoclaro}{HTML}{E4E6F3}
\definecolor{laranjaclaro}{HTML}{FDE8DC}

% ---- Aparência geral -----------------------------------------------------------
\usetheme{default}
\usefonttheme{professionalfonts}
\renewcommand{\familydefault}{\sfdefault}
\setbeamertemplate{navigation symbols}{}
\setbeamercolor{background canvas}{bg=fundo}
\setbeamercolor{normal text}{fg=black!85}
\setbeamercolor{structure}{fg=indigo}
\setbeamercolor{frametitle}{fg=indigo,bg=fundo}
\setbeamercolor{itemize item}{fg=laranja}
\setbeamercolor{itemize subitem}{fg=indigo}
\setbeamercolor{enumerate item}{fg=laranja}
\setbeamercolor{block title}{bg=indigo,fg=white}
\setbeamercolor{block body}{bg=indigoclaro,fg=black!85}
\setbeamercolor{alerted text}{fg=laranja}
\setbeamerfont{block title}{series=\bfseries,size=\normalsize}
\setbeamertemplate{blocks}[rounded][shadow=false]
\setbeamertemplate{itemize item}[circle]
\setbeamertemplate{itemize subitem}[circle]
\setbeamertemplate{enumerate item}{\insertenumlabel.}
\setbeamerfont{frametitle}{series=\bfseries,size=\Large}
\setbeamerfont{caption}{size=\scriptsize}

% título do slide: alinhado à esquerda com uma barra laranja
\setbeamertemplate{frametitle}{%
  \vspace{0.3cm}
  \begin{tikzpicture}[baseline=0pt]
    \fill[laranja] (0,-0.25) rectangle (0.14,0.55);
    \node[anchor=west,text=indigo,font=\Large\bfseries,text width=0.92\paperwidth] at (0.3,0.15) {\insertframetitle};
  \end{tikzpicture}
}

% rodapé: identificação e numeração
\setbeamertemplate{footline}{%
  \hbox to \paperwidth{\hspace*{0.6cm}%
    {\tiny\color{cinza}Técnicas de Avaliação de Políticas Públicas \,$\cdot$\, PPGGP EACH/USP \,$\cdot$\, Seminário Discente 2026}%
    \hfill {\tiny\color{cinza}\insertframenumber}\hspace*{0.6cm}}%
  \vspace*{0.25cm}%
}
\setlength{\leftmargini}{1.2em}

% ---- Macros --------------------------------------------------------------------
\newcommand{\fonte}[1]{\par\vfill{\tiny\color{cinza}Fonte: #1}}
\newcommand{\destaque}[1]{{\bfseries\color{laranja}#1}}
\newcommand{\ind}[1]{{\bfseries\color{indigo}#1}}
\newcommand{\cod}[1]{\texttt{\color{indigo}#1}}

% Slide de seção (fundo índigo)
\newcommand{\secao}[3]{%   #1 = número, #2 = título, #3 = tempo/subtítulo
  {\setbeamercolor{background canvas}{bg=indigo}%
  \begin{frame}[plain,noframenumbering]
    \begin{tikzpicture}[remember picture,overlay]
      \fill[laranja] (current page.south west) rectangle ++(0.5,\paperheight);
      \node[anchor=west,text=white,font=\Huge\bfseries,text width=0.8\paperwidth,align=left]
        at ([xshift=1.5cm,yshift=0.6cm]current page.west) {#2};
      \node[anchor=west,text=laranjaclaro,font=\large,align=left]
        at ([xshift=1.5cm,yshift=-1.2cm]current page.west) {#1 \,$\cdot$\, #3};
    \end{tikzpicture}
  \end{frame}}%
}

% Caixa "Hora do R" (demonstração em R)
\newenvironment{horaR}[1]{%
  \begin{beamercolorbox}[wd=\textwidth,rounded=true,leftskip=0.4cm,rightskip=0.4cm,sep=0.25cm]{laranjabox}
  {\bfseries\color{laranja}\large $\blacktriangleright$ Hora do R \,$\cdot$\, \cod{#1}}\par\smallskip
}{%
  \end{beamercolorbox}
}
\setbeamercolor{laranjabox}{bg=laranjaclaro,fg=black!85}

% Slide de atividade: bloco laranja
\newenvironment{atividadebox}[1]{%
  \setbeamercolor{block title}{bg=laranja,fg=white}
  \setbeamercolor{block body}{bg=laranjaclaro,fg=black!85}
  \begin{block}{#1}
}{%
  \end{block}
  \setbeamercolor{block title}{bg=indigo,fg=white}
  \setbeamercolor{block body}{bg=indigoclaro,fg=black!85}
}

% Trechos de código R
\lstdefinestyle{R}{
  language=R, basicstyle=\ttfamily\scriptsize, keywordstyle=\color{indigo}\bfseries,
  commentstyle=\color{verde}\itshape, stringstyle=\color{laranja},
  backgroundcolor=\color{white}, frame=single, rulecolor=\color{indigoclaro},
  breaklines=true, showstringspaces=false, columns=fullflexible, keepspaces=true,
  literate={á}{{\'a}}1 {à}{{\`a}}1 {â}{{\^a}}1 {ã}{{\~a}}1 {é}{{\'e}}1 {ê}{{\^e}}1
           {í}{{\'i}}1 {ó}{{\'o}}1 {ô}{{\^o}}1 {õ}{{\~o}}1 {ú}{{\'u}}1 {ç}{{\c c}}1
           {Á}{{\'A}}1 {É}{{\'E}}1 {Í}{{\'I}}1 {Ó}{{\'O}}1 {Ú}{{\'U}}1 {Ç}{{\c C}}1
           {°}{{\textdegree}}1
}
\lstset{style=R}

% Estilos TikZ comuns
\tikzset{
  cx/.style={draw=indigo,fill=white,rounded corners=3pt,thick,align=center,inner sep=5pt,font=\small},
  cxl/.style={draw=laranja,fill=laranjaclaro,rounded corners=3pt,thick,align=center,inner sep=5pt,font=\small},
  cxi/.style={draw=indigo,fill=indigoclaro,rounded corners=3pt,thick,align=center,inner sep=5pt,font=\small},
  seta/.style={-{Stealth[length=2.6mm]},thick,color=indigo},
  setal/.style={-{Stealth[length=2.6mm]},thick,color=laranja},
}
